% lw-vs English manuscript
% Build: pdflatex main -> bibtex main -> pdflatex main -> pdflatex main
\documentclass[lettersize,journal]{IEEEtran}

\usepackage{amsmath,amssymb}
\usepackage{booktabs}
\usepackage{graphicx}
\usepackage{array}
\usepackage{multirow}
\usepackage{url}
\usepackage[hidelinks]{hyperref}
\usepackage{enumitem}
\usepackage{algorithm}
\usepackage{algpseudocode}

\graphicspath{{../../figures/}{../figures/}}

% Tighten float spacing a little without altering IEEEtran's layout contract.
\setlength{\textfloatsep}{9pt plus 2pt minus 2pt}
\setlength{\floatsep}{8pt plus 2pt minus 2pt}
\setlength{\intextsep}{8pt plus 2pt minus 2pt}

\newcommand{\dblfig}[3]{%
  \begin{figure*}[t]\centering
    \includegraphics[width=\textwidth]{#1}%
    \caption{#2}\label{#3}
  \end{figure*}}
\newcommand{\onefig}[4]{%
  \begin{figure}[t]\centering
    \includegraphics[width=\columnwidth]{#1}%
    \caption{#2}\label{#3}\vspace{#4}
  \end{figure}}

\begin{document}

\title{First-Round Threshold Timing Controls Conditional\\ Decoding Failure in CLASSIC-Style QC-MDPC Bit-Flipping:\\ A Key-Conditional Study at Calibrated Operating Points}

\author{[Author names and affiliations to be completed before submission]%
\thanks{[Corresponding-author e-mail to be completed.]}}

\maketitle

\begin{abstract}

Reliable decoding is a correctness requirement for QC-MDPC-based post-quantum cryptography, yet its failure tail depends on both threshold scheduling and secret-key structure. Aggregate decoding-failure rate (DFR) cannot answer a more basic question: on one and the same instance, does the temporal configuration of a threshold change conditional risk? We define FRO-2, which raises by two the first-iteration threshold of the classic threshold bit-flipping decoder (CLASSIC), and evaluate it on paired $(H,e,s)$ instances with key-clustered inference. At $(r,d,t)=(4801,45,92)$ it lowers conditional DFR 4.15--5.47-fold across three conditions, every key-cluster interval excluding one. The mechanism is identified: incorrect first-round flips fall about $9.5\times$, correct flips only $2.7\times$. A round-controlled ablation separates where an offset is applied from whether it persists: the first round is a two-sided window ($+2$ helps 4.15-fold, $-1$ hurts eightfold), a one-shot round-2 offset helps 1.21-fold, and the same offset sustained destroys the decoder (0.017). The effect transfers without retuning and replicates at a larger code, but not under a capture-prone $\kappa>d$ condition, where a near-codeword-aware decoder is decisive. Fourteen times lower baseline DFR leaves the risk ratio at 3.71 [2.69, 5.57], stable rather than strengthening. At the BIKE Level-1 parameters $(12323,71)$, with the error weight raised into the measurable band, the window is two-sided there as well (4.48 against 0.162); at the production weight no failure is measurable, and the measurable signals are the round-one margin ($+2$ cuts false first-round flips sixfold) and a small convergence cost. A counterfactual model predicting other offsets from round-one statistics fails out of sample; we report that negative result.
\end{abstract}

\begin{IEEEkeywords}
post-quantum cryptography, QC-MDPC codes, bit-flipping decoding, decoding failure rate, first-round threshold, key-conditional risk.
\end{IEEEkeywords}

% ===========================================================================
\section{Introduction}\label{sec:intro}

QC-MDPC codes provide compact representations of sparse parity-check structures and form the decoding core of BIKE, a code-based key-encapsulation mechanism studied in the NIST post-quantum process \cite{misoczki2013mdpc,bike2024spec,alagic2025nist}. Their practical appeal is coupled to a difficult correctness problem: the iterative decoder is intentionally simple, but its residual failure probability occupies an extreme tail that cannot be measured directly. More importantly, that tail is key-dependent in ways that security analyses exploit: reaction, timing, and weak-key analyses have shown that decoding behaviour can reveal or amplify structural properties of the secret parity-check matrix \cite{guo2016key,eaton2018timing,wang2023weak,gellersen2025}. Whether that dependence becomes visible as per-key failure-rate spread at \emph{measurable} operating points is itself an empirical question, and in our primary pool it is not resolved---the per-key failure counts are consistent with homogeneous Poisson variation, a bound we report rather than assume away (Section~\ref{sec:discussion}). The scientific problem is therefore not only how often a decoder fails on average, but how its schedule changes conditional risk across keys and structured error regimes, whatever the residual spread turns out to be.

Most decoder studies optimize a global error-correction target or a mean DFR. Such summaries are necessary but incomplete when keys are the long-lived experimental units: thousands of error samples generated for one key are not thousands of independent key observations. A schedule can improve an aggregate mean while concentrating failures on a small key subset, and a comparison made on different error pools confounds the schedule with sample difficulty. Three design requirements follow: pair decoder configurations on the same $(H,e,s)$ instances, retain per-key outcomes, and conduct cross-key inference with the key as the resampling unit.

This study asks whether the earliest decoder decision can be stabilized with a deliberately small intervention. We call the decoder whose schedule is perturbed the \emph{carrier}; here it is CLASSIC, and in the negative control it is a fixed-threshold variant. For CLASSIC we define FRO-$k$ as a one-shot offset of $+k$ applied to the first-round threshold only (FRO-1, FRO-2 and FRO-3 are $+1$, $+2$ and $+3$); FRO-2 is the schedule studied throughout. The intervention contains no adaptive state, no learned model, and no additional randomness. Its rationale is temporal rather than global: the initial syndrome is the heaviest, and a conservative first decision can prevent damaging flips before the later correction dynamics take over. Whether a gain is genuinely tied to the first-round state, or merely to the presence of an offset somewhere in the schedule, is what the round-controlled ablation below separates.

We make four contributions. First, we establish a paired, key-conditional evaluation of threshold interventions in which the key, rather than the individual error draw, is the resampling unit. Second, using that protocol we show that the safety of a threshold offset is governed by two independent properties of its temporal configuration --- which round it touches, and whether it is a one-shot perturbation or a sustained displacement: the first round is a sharp two-sided window, and sustaining an offset across later rounds is destructive on its own. Third, we connect the observed DFR reduction to a local unsatisfied-parity-check (UPC) tail model, which reveals an asymmetry --- incorrect flips are suppressed roughly $9.5\times$ whereas correct flips only about $2.7\times$ --- and we map the transfer of the effect across operating points, decoder semantics, iteration budgets, key families and a capture-prone structured-error boundary, positioning the intervention against modern QC-MDPC decoders. Fourth, we carry the measurement to the production code parameters. At $(r,d)=(12323,71)$, the exact Level-1 blocks, the error weight must be raised by about $28\%$ before the baseline failure rate enters the measurable band; there the first-round window is two-sided (risk ratios 4.48 and 0.162), while at the production error weight itself no failure is measurable and the measurable production signals are the round-one decision margin and a convergence cost --- both quantified here, together with an ablation of the first-round margin in the specification's own decoder. Alongside these we report a negative result: a counterfactual model that predicts other offsets' risk ratios from round-one statistics does not transport across offsets, for a reason we diagnose.

Section~\ref{sec:background} gives background, related work, and the boundaries of our claims. Section~\ref{sec:model} develops the local model and its three observable signatures. Section~\ref{sec:method} describes the experimental design and statistical protocol. Section~\ref{sec:results} reports results, and Section~\ref{sec:discussion} discusses scope and limitations.

% ===========================================================================
\section{Background and Related Work}\label{sec:background}

\subsection{QC-MDPC bit-flipping decoding}

Let $H=[H_0\mid H_1]$ be a binary parity-check matrix composed of two $r\times r$ circulant blocks with block weight $d$. For an error vector $e$ of weight $t$, decoding begins from the syndrome $s=He^{\top}$. A parallel bit-flipping decoder computes an unsatisfied-parity-check count $u_j$ for every candidate position and flips positions satisfying a threshold rule. The syndrome is then updated and the process repeats until it vanishes or the iteration budget is exhausted.

The carrier throughout this paper is the \emph{classic threshold bit-flipping decoder}: the parallel bit-flipping decoder just described, whose threshold is derived from a binomial model of the current syndrome weight \cite{misoczki2013mdpc,vasseur2021thesis}. Following the reference implementation we refer to it as CLASSIC, and we denote its unmodified schedule by $C0$; it is the paired control throughout. All decoders discussed below --- CLASSIC, BGF, BIKE-Flip, Backflip and BF-Max --- are members of this same parallel bit-flipping family and differ in how the threshold is scheduled. Decoder behaviour depends strongly on that schedule. A low threshold flips more correct error positions but can also introduce incorrect flips; a high threshold protects against false flips but can leave too few positions corrected, producing stagnation. BGF supplements an affine threshold with black and gray passes \cite{drucker2020shades}. BIKE-Flip is the tuned production decoder of BIKE specification v5.2 \cite{bike2024spec}. In the specification's notation its threshold is $\mathrm{threshold}(S,i,S_0)=\max(f_t(S),T_i(S_0))$, where $f_t$ is affine in the current syndrome weight $S$ and the floors $T_i(S_0)$ carry a fixed margin $\Delta$ in \emph{every} iteration: $T_1=f_t(S_0)+\Delta$; $T_2$ and $T_3$ blend $f_t(S_0)$ with the decoder's floor $(v{+}1)/2$ and add $\Delta$ as well; $T_i=(v{+}1)/2+\Delta$ for $i\ge4$. At Level-1 $\Delta=3$ with seven fixed iterations. That margin is a different quantity from the offset $\delta$ studied here, and it is not confined to the first round. The published rationale is that failures originate in iterations two and three, when a few correct positions carry extreme counters while the syndrome-derived threshold has already fallen \cite{sendrier2024seminar}. We therefore stress that FRO-2 is \emph{not} the specification's schedule: it is a one-shot $+2$ on the first iteration of the CLASSIC schedule, whereas the specification's margin is a permanent part of a differently shaped one. Backflip stores limited state and permits time-to-live-based re-flipping \cite{sendrier2020lowdfr}. BF-Max selects a maximum-counter position per iteration with random tie-breaking \cite{baldelli2025bfmax}. Recent near-codeword-aware methods address error-floor trajectories through more targeted state and structure \cite{baldelli2026ncw,santini2026flip}.

FRO-2 is not proposed as a new decoder family. It is a one-parameter perturbation of CLASSIC designed to isolate the causal value of first-round conservatism. That narrow scope is what makes it possible to distinguish a temporal threshold effect from general schedule tuning.

\subsection{Conditional failure and weak structures}

The failure probability of a QC-MDPC decoder is conditional on the secret key. Guo, Johansson, and Stankovski showed that decoding errors expose distance-spectrum information \cite{guo2016key}, and subsequent timing and weak-key studies sharpened the attack surface \cite{eaton2018timing,wang2023weak,gellersen2025}. Nosouhi \emph{et al.} implemented BIKE, quantified weak-key effects on average DFR, and proposed a key-check mechanism; that line of work changes the accepted-key population \cite{nosouhi2023weakkey}. Structural key filters can reduce exposure at key generation \cite{matthews2025combinatorial}, while decoder countermeasures attempt to make difficult structures less influential. These approaches are complementary: filtering changes the population of accepted keys; scheduling changes behaviour conditional on a key.

The phrase ``weak key'' covers distinct mechanisms and must not be used as a single undifferentiated category. Our Type-I construction forces a run of 12 consecutive ones in one parity-check block. At the tested dimensions this is an adversarially rare event, used to stress the decoder rather than to estimate population-average risk. The shallow near-codeword condition injects errors that intersect a selected near-codeword in only $\kappa=12$ positions, with $\kappa<d=45$; because the resulting distance to that near-codeword remains above the target error weight, capture is excluded by construction. It is therefore a \emph{shallow-intersection control}, not an error-floor or gathering proxy. A separate fixed-key-batch experiment varies $\kappa$ beyond $d$ as a capture-prone structured-error stress test; the observed failures there do not by themselves establish entry into a dynamical capture basin.

\subsection{Boundaries of the contribution}

The positioning of this work needs to be explicit, because it bounds every claim below.

\begin{itemize}[leftmargin=1.3em,itemsep=0.12em,topsep=0.2em]
\item This is a \emph{calibrated decoder-mechanism study}. It is not a DFR theorem at production parameters, not a weak-key recovery attack, and not an IND-CCA security proof.
\item The main results sit at $10^{-4}$--$10^{-3}$ calibrated pressure points, far above BIKE's cryptographic target. No extrapolation of the failure-rate \emph{level} is attempted: the production-parameter section reports a measured operating point at the production blocks with a raised error weight, and mechanism-level statistics at the production error weight itself.
\item Comparisons with BIKE-Flip and Backflip are references under unequal semantics and tuning, not dominance claims.
\item Conditions labelled ``capture-prone'' are structured-error stress conditions; we record neither terminal trajectory distances nor basin fractions.
\item Threshold randomization is rejected within the tested design space. This does not prove that every randomized decoder is ineffective---in-place and stateful randomized methods differ substantially \cite{baldi2021inplace}.
\item Round-one statistics are a diagnostic, not a validated predictor: the per-instance model developed in Section~\ref{sec:counterfactual} fails out of sample, and Section~\ref{sec:discussion} gives the reason.
\item The production-parameter point is the production \emph{code}, not the production \emph{operating point}: only the error weight is raised, and the measured effect there is not evidence about the failure rate at the designed weight.
\end{itemize}

\subsection{Related modelling work}

Two lines of prior work overlap with this mechanism study and delimit what is new here. The first models decoding failure from round-level state. Kim and Park predict LDPC decoding failure with a logistic model over iteration index and unsatisfied-check count for early termination \cite{kim2021logistic}; Paiva and Terada use the average number of uncorrected errors remaining after the first iteration as a small-sample DFR proxy across code sizes \cite{paiva2022pickyfix}; Santini \emph{et al.} build a code-specific conservative model of the failure rate from the probability that the first round fails to reduce the error count \cite{santini2019conservative}; and Annechini \emph{et al.} estimate failure rates for $(v,w)$-regular codes from the syndrome-weight distribution \cite{annechini2024estimation}. The second line models the failure rate itself: Sendrier and Vasseur's Markov chain over syndrome weight and error weight \cite{sendrier2019dfr}, Tillich's asymptotic law \cite{tillich2018dfp}, Vasseur's two-point extrapolation framework \cite{vasseur2021thesis}, key-dependent Markov models for the error floor \cite{arpin2025errorfloor}, and a closed-form model for the exact three-iteration parallel decoder \cite{annechini2025closedform}, which explicitly warns that extrapolating from a proxy decoder is hazardous. None of this prior work uses round-one statistics \emph{counterfactually}---to predict the risk ratio of a different threshold schedule---and our own attempt to do so fails out of sample for a diagnosable reason (Section~\ref{sec:discussion}); we report that attempt and its failure below rather than presenting round-one statistics as predictive.

% ===========================================================================
\section{A Local Model for the First-Round Window}\label{sec:model}

Let $U_{i,j}$ denote the UPC count for position $j$ in iteration $i$. Conditional on syndrome weight and local correlations, a first-order model is
\begin{equation}
\begin{aligned}
U_{i,j}\mid j\in\mathrm{err} &\sim \mathrm{Bin}(d,p_{1,i}),\\
U_{i,j}\mid j\notin\mathrm{err} &\sim \mathrm{Bin}(d,p_{0,i}),
\end{aligned}
\end{equation}
with $p_{1,i}>p_{0,i}$. For a threshold $\tau_i$ define the correct- and incorrect-flip tails $q_{1,i}(\tau_i)$ and $q_{0,i}(\tau_i)$. FRO-2 replaces only $\tau_1$ by $\tau_1+2$, so the two tails are multiplied by
\begin{equation}
A_1=\frac{q_{1,1}(\tau_1+2)}{q_{1,1}(\tau_1)},\qquad
A_0=\frac{q_{0,1}(\tau_1+2)}{q_{0,1}(\tau_1)} .
\end{equation}

When the threshold lies in the steep incorrect-flip tail but retains mass for true errors, the model predicts $A_0\ll A_1$: false flips are suppressed more strongly than correct flips. The predicted trajectory is a larger round-one residual but a lower probability of entering a cascade or stagnation basin.

The approximation conditions on syndrome weight and neglects inter-UPC dependence; it is used to predict \emph{relative first-round changes}, not asymptotic DFR. It yields one signature directly and frames two further hypotheses that we test separately. Directly:

\begin{enumerate}[label=(\roman*),leftmargin=1.7em,itemsep=0.1em,topsep=0.2em]
\item incorrect first-round flips fall faster than correct first-round flips;
\item a \emph{sustained} $+2$ offset need not help, even when it avoids the first round entirely;
\item a small offset need not help under a capture-prone $\kappa>d$ structured-error condition.
\end{enumerate}

The remaining two claims --- that a \emph{sustained} offset is harmful even when it avoids round~1, and that a small offset need not help under a capture-prone $\kappa>d$ condition --- lie outside what this first-round model can predict; the model says nothing about the sign of the offset or about later rounds, and we treat them as hypotheses tested by experiment rather than as model consequences. The mechanism experiment, the temporal ablation with its round controls, and the high-overlap stress experiment in Section~\ref{sec:results} test them.

% ===========================================================================
\section{Experimental Design}\label{sec:method}

\subsection{Evidence tiers and decision structure}

The pipeline comprises calibration, candidate screening, held-out confirmation, formal measurement, and preregistered sensitivity analyses. Fig.~\ref{fig:design} summarises the protocol and the evidence structure, and because the reader must be able to judge the strength of each claim, Table~\ref{tab:tiers} states the scale and the inferential authority of each of the four evidence tiers used in this paper. Every quoted number is reproducible from frozen seeds and configurations; tier-3 and tier-4 results are labelled exploratory or prototype throughout and never carry a primary conclusion.

\begin{table}[t]
\centering
\caption{Evidence tiers, by the scale that dominates each. Individual experiments are labelled in the text; tier-3 results are marked exploratory wherever they appear and do not carry primary conclusions.}
\label{tab:tiers}
\footnotesize
\fontsize{6.2}{7.4}\selectfont
\setlength{\tabcolsep}{3pt}
\begin{tabular}{@{}clll@{}}
\toprule
Tier & Role & Scale & Authority \\
\midrule
1 & formal primary and its controls & 92--94 keys $\times$ 3{,}000; held-out 48 $\times$ 1{,}000 & confirmatory \\
2 & preregistered confirmation & 50 keys $\times$ 1{,}000 & confirmatory \\
3 & exploratory ablations & 15--30 keys $\times$ 1{,}000--2{,}000 & exploratory \\
4 & prototype engineering & 30 keys $\times$ 500 & feasibility \\
\bottomrule
\end{tabular}
\end{table}

Whole keys, not error instances, were assigned across search, validation, and test stages. Search evaluated nine schedules: the unmodified carrier, uniform offsets, first-round offsets, a two-round variant, and coin-selected threshold variants. Successive halving retained candidates using the worst condition-group one-sided 95\% DFR upper bound, subject to a random-group non-degradation constraint. The selected schedule, FRO-2, raises only the first-round threshold by two. No later threshold differs from C0.

\dblfig{fig1_study_design.pdf}
{Study design. (a) Paired decoding protocol: C0 and FRO-2 process the same $(H,e,s)$, and the true error is used only to generate the instance and for instrumented diagnostics. (b) Threshold timing: FRO-2 raises the threshold by two in round~1 only and is identical to C0 thereafter. (c) Scale and role of the four evidence tiers; tier-3 and tier-4 results are labelled exploratory or prototype throughout.}
{fig:design}

\subsection{Operating points and condition groups}

The primary measurable point is $P_1'=(r,d,t)=(4801,45,92)$ with two circulant blocks --- the measurable variant of the screening point $P_1=(4801,45,84)$ used during candidate search (Section~\ref{sec:negative}). Its block length and block weight are those of the parameter set Misoczki \emph{et al.} recommended for 80-bit security, $(n,r,\omega,t)=(9602,4801,90,84)$ \cite{misoczki2013mdpc}, whose per-block weight is $d=\omega/2=45$. We retain $(r,d)$ and raise the error weight from the recommended $84$ to $92$ so that the control decoder fails often enough for direct simulation to resolve. A calibration sweep over $t\in\{86,88,90,92,95\}$ confirmed the placement: at $t\le90$ no C0 failure was observed in 1{,}600 calibration draws, and at $t=92$ the calibration estimate was $3.1\times10^{-3}$. The point is therefore a published recommended parameter set deliberately over-driven for measurability, not an arbitrary surrogate; it remains far from BIKE's production parameters, and no proportional scaling to NIST security levels is implied.

Three formal conditions were used:

\begin{itemize}[leftmargin=1.3em,itemsep=0.12em,topsep=0.2em]
\item \textbf{RAND}: independently generated random QC-MDPC keys with uniformly sampled weight-92 errors.
\item \textbf{TYPE-I}: keys constrained to contain a 12-position consecutive run in one block, paired with uniformly sampled errors.
\item \textbf{NCW-shallow}: keys admitted by a near-codeword criterion, with errors constructed to intersect the selected near-codeword in 12 positions. Because $12<d$, this is a shallow-intersection control, not a capture condition.
\end{itemize}

The formal pool contains 92 RAND, 94 TYPE-I, and 92 NCW-shallow keys, each receiving 3{,}000 errors, giving 834{,}000 observations per arm. A pure-fresh subset contains 212 keys not used in any interim decision, totalling 636{,}000 observations per arm; that subset is the preregistered gate criterion, and the full-pool figure is reported as a sensitivity result.

The transfer point is $P_2=(3001,37,66)$, calibrated independently so that C0 again had measurable failures. The offset remained exactly $+2$ in round~1 and was not retuned. Thirty keys per condition and 3{,}000 errors per key gave 90{,}000 observations per arm and condition.

BIKE Level-1 uses the native $(r,d,t)=(12323,71,134)$ parameter set and is retained as a production-boundary hold-out, not as a primary measurable result. Two production-scale checks were performed. First, offsets $0$, $+2$, and $+4$ were each run on one million instances in RAND, TYPE-I, and near-codeword conditions, for nine million total runs. Second, a specification-exact BIKE-Flip implementation with official Level-1 constants was compared with C0 and FRO-2 on 16 keys and 500 errors per condition.

The preregistered expanded confirmation used an independent master seed (\texttt{20260912}), with 50 regenerated keys per condition and 1{,}000 errors per key at $P_2$ and at the supplemental calibrated point $P_3=(4001,41,82)$. The six schedules shared 300{,}000 unique $(H,e,s)$ inputs, yielding 1{,}800{,}000 decoder-arm runs. This experiment is separate from the seed-\texttt{20260911} exploratory matrix and was used to confirm the first-round placement effect and the condition dependence of offset magnitude.

\subsection{Decoder semantics and cross-implementation validation}

The Python implementation represents arithmetic in $\mathbb{F}_2[x]/(x^r-1)$ and includes reference and accelerated UPC paths. Unit tests compare syndrome construction, UPC values, flip semantics, generators, failure classification, and statistical routines against direct definitions. BIKE-Flip's affine coefficients and integer thresholds were reconstructed from specification v5.2 and validated against the published Level-1, Level-3, and Level-5 schedules. BF-Max uses pseudorandom tie-breaking among positions with equal maximum counters, matching the published algorithm.

Cross-implementation validation was assessed at two levels. First, a separately coded dynamic-distance Python path observed 11/20{,}000 failures versus 136/200{,}000 in the C reference, giving $z=-0.68$, within the predeclared consistency band. Second, the C reference evaluated one million instances per condition and reproduced the direction of FRO-2 on RAND and Type-I (1.59-fold and 2.01-fold), while the NCW comparison was unresolved at that sample size. Both implementations were developed within the same research project, so this checks implementation agreement rather than independent-laboratory reproducibility. The dynamic threshold also uses simulation-privileged true-error distance.

The main study uses a static target error weight $t$ in the threshold calculation and declares success at zero syndrome. In a separate exact-recovery audit, we tracked the parity of all decoder flips and compared the resulting estimated error with the generated error; the decoder itself did not receive the true error. The audit covered all 278 formal $P_1'$ keys and 834{,}000 decodings per arm. It found 795 exact-recovery failures for C0 and 167 for FRO-2, matching the zero-syndrome failure totals, with no zero-syndrome miscorrection observed in either arm. This is empirical agreement on the audited samples, not a logical equivalence between zero syndrome and recovery of the generated error.

Algorithm~\ref{alg:paired} states the paired procedure. C0 and FRO-2 process the same $(H,e,s)$; the true error is used only to generate the paired instance, to label correct versus incorrect flips in the diagnostic analysis, and to compare estimated errors in the separate recovery audit.

\begin{algorithm}[t]
\caption{Paired C0/FRO-2 decoding and evaluation}
\label{alg:paired}
\footnotesize
\begin{algorithmic}[1]
\Require key $H$, error $e$, threshold rule $\tau(i,s)$, budget $T$
\State $s_0 \gets H e^{\top} \bmod 2$ \Comment{experiment driver}
\For{$\mathrm{arm}\in\{\mathrm{C0},\mathrm{FRO\text{-}2}\}$}
  \State $s\gets s_0$
  \For{$i\gets 1,\dots,T$}
    \State $\delta \gets 2$ if $(\mathrm{arm}=\mathrm{FRO\text{-}2}$ and $i=1)$ else $0$
    \State $u \gets \mathrm{UPC}(H,s)$
    \State $F \gets \{j: u_j \ge \tau(i,s)+\delta\}$
    \State $s \gets s \oplus H\,\mathbf{1}_F^{\top} \bmod 2$
    \State record flips and syndrome weight
    \If{$s=0$} record success and break \EndIf
  \EndFor
  \If{$s\neq 0$} record failure \EndIf
\EndFor
\State \Return the paired C0/FRO-2 outcomes
\end{algorithmic}
\end{algorithm}

For a fixed UPC routine and iteration budget, FRO-2 has the same asymptotic work and memory as C0: $O(Trd)$ for the sparse QC-MDPC update and $O(r)$ decoder state. Its only algorithmic addition is one integer threshold offset in round~1; it adds no UPC pass, no persistent state, and no random draw. The measured cost difference therefore reflects the changed convergence trajectory and the implementation, not a new asymptotic operation.

\subsection{Outcomes and statistical analysis}

The primary outcome is the arm DFR within each condition group, the number of failed decodings divided by the number attempted. We call the design \emph{key-conditional} in a separate sense: inference is made at the key level, with outcomes paired per instance and keys resampled, so the conditioning that defines the statistic and the unit of inference must not be conflated. Individual outcomes remain paired because C0 and FRO-2 process the same $(H,e,s)$. Secondary outcomes include the one-sided 95\% Wilson upper bound, per-key failure counts, convergence iterations, and a failure taxonomy.

The principal cross-key effect is the risk ratio $\mathrm{RR}=\widehat{\mathrm{DFR}}(\mathrm{C0})/\widehat{\mathrm{DFR}}(\mathrm{FRO\text{-}2})$. Its 95\% interval is estimated by a nonparametric cluster bootstrap with keys as clusters, 10{,}000 resamples, and seed \texttt{20260905}; each resample draws keys with replacement and retains all errors for each selected key. We also report absolute DFR differences, exact two-sided sign tests on per-key paired failure counts, and exact McNemar tests on paired instance outcomes. Sign-test and McNemar probabilities are Holm-adjusted across the three formal conditions. A beta-binomial sensitivity model estimates a schedule odds ratio and an intra-key overdispersion term; at the resolution of these data the fitted overdispersion is indistinguishable from zero ($p\approx0.49$), so the model collapses to the pooled risk ratio and returns a narrower interval than the key-cluster bootstrap. We therefore report it only as a directional sensitivity check, never as an independent precision statement.

This paper reports on the order of fifty cluster-bootstrap intervals across its results sections. Holm adjustment covers the sign and McNemar tests of the three formal conditions, but the risk-ratio intervals are marginal by construction: each is a 95\% interval for a single contrast, not a simultaneous coverage statement for the family, and they are not adjusted. Phrases such as ``every interval excludes one'' should be read per-contrast, not as a joint 95\% claim; the one place where a multi-cell claim is made, Table~\ref{tab:temporal}, says so explicitly.

For nonzero counts, DFR is reported with a one-sided 95\% Wilson upper bound; for zero failures we give the rule-of-three bound $3/N$ and never interpret the estimate as proof of zero risk. A single zero-count convention is used throughout, with no implicit continuity correction.

The preregistered operational gate compares the worst condition-group one-sided upper bounds on a base-10 log scale and requires at least 0.5 $\log_{10}$ improvement without random-group degradation. This is a decision rule, not a formal non-inferiority hypothesis test. The originally proposed heterogeneity criterion---a P90--P10 per-key DFR span of at least 3 $\log_{10}$ units---was not met; observed spans were 0.40--0.48 $\log_{10}$ units.

\subsection{Implementation and reproducibility}

Keys, error pools, candidate definitions, thresholds, and random choices derive from master seed \texttt{20260901} through recorded SHA-256 derivations. Raw per-key counts, configuration files, analysis outputs, and the C build shims are archived. Figure source tables are generated directly from the archived JSON by \path{lw-vs/scripts/export_data.py}, so every number cited in the text is traceable to a specific archive file.

The formal static runs used Python 3.12, NumPy 2.4.6, and a single process; the C reference used Zig-CC at \texttt{-O2} without AVX. The archive does not record the CPU model or RAM, so wall-clock values indicate computational scale only and are not suitable for cross-machine performance comparisons.

% ===========================================================================
\section{Results}\label{sec:results}

\subsection{Temporal structure: a two-sided first round and the cost of sustained offsets}\label{sec:temporal}

The effect of a threshold offset depends on two distinct properties of its temporal placement: which round it touches, and whether it is a one-shot perturbation or a sustained displacement of the schedule. The earlier design could not separate them, because the uniform offset differs from a first-round offset in the later rounds as well as in round~1. We therefore added the missing arms: a sustained offset from round~2 onwards, which never touches round~1, and a first-round offset of the opposite sign. These two arms are \emph{post hoc} relative to the original design --- they were specified after the confound in the uniform-versus-first-round contrast became apparent, and they are not preregistered. They are reported at full formal scale on all three conditions, but they come from a single operating point, and we flag that limitation rather than present them as preregistered confirmations. All four configurations were measured on the full formal pool at $P_1'$, with C0 and FRO-2 reproducing the archived counts exactly as a pipeline check (Fig.\ref{fig:temporal}d).

The first round is a sharp \emph{two-sided} window. Adding two to its threshold reduces DFR from $1.022\times10^{-3}$ to $2.464\times10^{-4}$ (risk ratio 4.15 [3.29, 5.38], 79 of 92 keys improving). Subtracting one from the same round is comparably destructive in the opposite direction: DFR rises to $8.10\times10^{-3}$, a risk ratio of 0.126 [0.111, 0.142], and \emph{every one of the 92 keys} is worse. A conservative first decision helps and a permissive one hurts, by similar magnitudes and with no key-level exceptions; the direction of the offset matters as much as its location.

The second property is independent of the first, and on this pool the two can be separated directly. Applying $+2$ from round~2 onwards, leaving round~1 untouched, raises DFR to $6.08\times10^{-2}$ (risk ratio 0.0168 [0.0149, 0.0187]), with every key worse. Applying $+2$ to round~2 \emph{alone}, on the same keys and the same error pools, instead \emph{lowers} DFR to $8.48\times10^{-4}$ (risk ratio 1.21 [1.07, 1.36]). The two configurations differ only in whether the offset persists past round~2, and they fall on opposite sides of unity in every condition: 1.21--1.45 when one-shot, 0.012--0.021 when sustained. Because both arms share the pool, this contrast does not carry the operating-point confound that the uniform-versus-first-round comparison did. Both new arms reproduce in the other two conditions: the first-round negative offset gives risk ratios of 0.109 [0.099, 0.120] in TYPE-I and 0.101 [0.085, 0.120] in NCW, and the sustained later-round offset gives 0.0210 [0.0192, 0.0229] and 0.0118 [0.0100, 0.0137], with no key improving in any of the six cells. The destructive configuration is therefore a \emph{sustained} displacement of the later-round schedule, not the mere presence of an offset outside round~1.

The remaining panels place this against the broader landscape. Across three conditions and two operating points, first-round offsets of $+1$, $+2$ and $+3$ all lift the risk ratio into the range 1.76--4.62 with all six cell intervals excluding one, while a uniform offset collapses it to 0.009--0.019 (Fig.\ref{fig:temporal}a--c). The two readings are consistent: every one-shot configuration tested improves the decoder --- a first-round offset by a large factor and a round-2-only offset by a small one --- whereas every sustained configuration is destructive regardless of where it starts.

\begin{table*}[t]
\centering
\caption{Temporal ablation: the fate of one $+2$ offset applied in different rounds (preregistered expanded confirmation, 50 keys $\times$ 1{,}000 errors per key). The risk ratio is C0 DFR divided by the schedule DFR, so values above one favour the schedule; bold marks the best first-round offset in each row, showing that FRO-2 is not optimal in every row.}
\label{tab:temporal}
\footnotesize
\begin{tabular}{@{}llcccccc@{}}
\toprule
\multirow{2}{*}{Point} & \multirow{2}{*}{Condition} & \multirow{2}{*}{C0 DFR}
 & \multicolumn{3}{c}{first-round offset} & \multirow{2}{*}{round-2 only} & \multirow{2}{*}{uniform $+2$} \\
\cmidrule(lr){4-6}
 & & & FRO-1 & FRO-2 & FRO-3 & & \\
\midrule
$P_2$ & RAND        & $2.56\times10^{-3}$ & \textbf{2.56} & 2.37 & 1.91 & 1.13 & $0.013$ \\
$P_2$ & TYPE-I      & $3.14\times10^{-3}$ & 2.31 & \textbf{2.85} & 2.15 & 1.29 & $0.019$ \\
$P_2$ & NCW-shallow & $1.20\times10^{-3}$ & 1.76 & \textbf{3.00} & 2.50 & 1.36 & $0.009$ \\
$P_3$ & RAND        & $1.10\times10^{-3}$ & 2.29 & \textbf{2.62} & \textbf{2.62} & 1.34 & $0.010$ \\
$P_3$ & TYPE-I      & $1.70\times10^{-3}$ & \textbf{3.27} & \textbf{3.27} & 2.18 & 1.31 & $0.016$ \\
$P_3$ & NCW-shallow & $7.40\times10^{-4}$ & 3.70 & \textbf{4.62} & 3.08 & 1.54 & $0.009$ \\
\bottomrule
\end{tabular}
\end{table*}

This is mechanistically informative. If the benefit merely came from moving a calibrated operating point away from a threshold cliff, applying $+2$ uniformly should improve the decoder; its large degradation rules that out. A uniform offset differs from a first-round offset in two ways at once --- it touches round~1, and it displaces every later round --- and the round controls above show that the second of these is sufficient on its own to destroy the decoder. Table \ref{tab:temporal} gives point estimates for all six condition cells: the selected FRO-2 lies between 2.37 and 4.62 while the uniform offset lies between 0.009 and 0.019, and FRO-1 and FRO-3 span a similar range, indicating that the \emph{first-round window itself} rather than one particular magnitude is what matters.

Notably, the round-2-only offset is not inert: its interval excludes one in four of the six cells (both TYPE-I and both NCW cells), but the magnitude is only about 1.1--1.5-fold, far below the first-round offsets. The seed-\texttt{20260911} exploratory matrix (12 keys $\times$ 400) observed only a 1.00--1.33 signal, too underpowered to resolve it; this explains why an earlier analysis read it as showing no stable benefit. We take the 50-key $\times$ 1{,}000 preregistered confirmation as authoritative here.

These are six marginal intervals, not an unadjusted simultaneous 95\% coverage statement, and the best offset magnitude also depends on operating point and key condition---at $P_3$/RAND both FRO-1 and FRO-3 slightly outperform FRO-2---so FRO-2 is a \emph{validated representative intervention}, not a universal optimum.

\dblfig{fig4_temporal_dissociation.pdf}
{Temporal dissociation. (a)--(c) Risk ratio against C0 for each schedule in the three conditions; the dashed line is the no-effect reference $\mathrm{RR}=1$. Filled circles are $P_2$ and squares are $P_3$. First-round offsets (FRO-1/2/3) sit well above one (1.76--4.62), the round-2-only offset stays close to but slightly above the reference (1.13--1.54), and the uniform $+2$ offset collapses to the $10^{-2}$ range. (d) Round controls on the full formal pool at $P_1'$, all decoded on the same keys and error pools. The first round is two-sided: $r_1$+2 helps 4.15-fold and the first-round $-$1 arm hurts 8-fold, and the one-shot versus sustained distinction is resolved within the pool: $+2$ at round 2 alone lowers DFR (risk ratio 1.21), whereas $+2$ from round 2 onwards destroys the decoder (0.0168) even though round 1 is untouched. Error bars are 95\% key-cluster bootstrap intervals in (a)--(c); labels in (d) are risk ratios against C0.}
{fig:temporal}

\subsection{Primary formal comparison}\label{sec:formal}

FRO-2 reduces failures in all three formal conditions (Table~\ref{tab:main}, Fig.~\ref{fig:main}). In RAND, failures fall from 282/276{,}000 to 68/276{,}000, corresponding to DFRs of $1.022\times10^{-3}$ and $2.464\times10^{-4}$ (4.15-fold). In TYPE-I they fall from 372/282{,}000 to 68/282{,}000, i.e. $1.319\times10^{-3}$ to $2.411\times10^{-4}$ (5.47-fold); in NCW from 141/276{,}000 to 31/276{,}000, i.e. $5.109\times10^{-4}$ to $1.123\times10^{-4}$ (4.55-fold).

\begin{table}[t]
\centering
\caption{Primary formal results at $(4801,45,92)$. RR is C0 DFR divided by FRO-2 DFR; intervals are 95\% key-cluster bootstrap intervals and $p$ is the Holm-adjusted sign test.}
\label{tab:main}
\footnotesize
\setlength{\tabcolsep}{3pt}
\begin{tabular}{@{}lcccc@{}}
\toprule
Condition & C0 fail/trials & FRO-2 fail/trials & RR [95\% CI] & Holm $p$ \\
\midrule
RAND & 282/276{,}000 & 68/276{,}000 & 4.15 [3.29, 5.38] & $4.9\times10^{-17}$ \\
TYPE-I & 372/282{,}000 & 68/282{,}000 & 5.47 [4.25, 7.28] & $3.0\times10^{-22}$ \\
NCW & 141/276{,}000 & 31/276{,}000 & 4.55 [3.21, 6.82] & $7.7\times10^{-12}$ \\
\bottomrule
\end{tabular}
\end{table}

All cluster intervals exclude 1, and the corresponding absolute DFR-difference intervals exclude 0. In per-key counts, FRO-2 is better on 233 keys, worse on 18, and tied on 27 (of 278); exact McNemar tests on paired instances are also significant after Holm correction. In the preregistered pure-fresh holdout, the worst-group log-upper-bound improvement is 0.687, exceeding the 0.5 operational gate, and RAND satisfies the non-degradation constraint. For comparison the full formal pool has a worst-group improvement of 0.679. The two agree in direction, but \emph{the pure-fresh value of 0.687 is the gate criterion}; the full-pool value is a sensitivity result, because the formal pool extends an interim pool rather than forming a wholly independent one.

The pure-fresh subset gives risk ratios of 4.69 [3.58, 6.34], 5.52 [4.24, 7.52], and 3.69 [2.53, 5.78]. A beta-binomial sensitivity model yields FRO-2/C0 odds ratios of 0.241, 0.183, and 0.220 in the full pool, all intervals excluding 1; the fitted extra key-level overdispersion is near zero at the resolution of these data. Together these analyses show that the aggregate result is produced neither by a handful of keys nor by interim reuse.

The exact-recovery audit rules out the zero-syndrome miscorrection concern: C0's 795 and FRO-2's 167 exact-recovery failures correspond one-for-one with the zero-syndrome failure totals, and the zero-syndrome miscorrection count is 0 in both arms.

\dblfig{fig2_main_results.pdf}
{Primary formal results. (a) Conditional DFR with one-sided 95\% upper caps in the three conditions. (b) Risk ratio C0/FRO-2; horizontal lines are 95\% key-cluster bootstrap intervals and the dashed line is the no-effect reference. (c) Paired per-key failure counts; the dashed line is the identity (no difference) and points are keys, not independent error samples.}
{fig:main}

\subsection{Mechanism: asymmetric suppression of incorrect flips}\label{sec:mechanism}

Direct instrumentation reveals a consistent first-round trade-off (Fig.~\ref{fig:mechanism}a). C0 averages about 29 flips in round~1, of which about 8 are incorrect. FRO-2 averages about 9 flips, of which about 0.8 are incorrect. \emph{Incorrect} first-round flips therefore decline by roughly $9.5\times$, whereas \emph{correct} first-round flips decline only about $2.7\times$. Because FRO-2 also postpones many correct flips, its mean residual error weight after round~1 is about 6 positions higher. Nevertheless, final failures in the instrumented sample fall from 21 to 7 in RAND, 25 to 6 in TYPE-I, and 10 to 1 in NCW, with mean iterations increasing by about 0.25--0.28.

The interpretation is that FRO-2 does not correct more of the error immediately; it defers more true errors while injecting far fewer false ones. Later rounds then recover the deferred corrections from a cleaner support configuration. On the subset rescued by FRO-2, the baseline makes approximately 16 incorrect first-round flips whereas FRO-2 makes 1.4--2.2.

These measurements use the true error support and are therefore \emph{diagnostic instrumentation}, not information available to a deployed decoder. That distinction matters: the analysis can explain why a schedule succeeds without implying a new implementable decision signal.

Failure-mode composition (Fig.~\ref{fig:mechanism}b) reports stagnation as the largest category: the mode counts for C0 failures are 198/282, 257/372, and 101/141 in RAND, TYPE-I, and NCW respectively. The archived counts were produced by a classifier whose early-misflip window was $\lvert\mathrm{flips}\rvert/4$ rather than the preregistered $\lceil T/4\rceil$, so we reran the classification on the full formal population under the preregistered definition. Every one of the 30 mode cells is identical, and the per-cell comparison ships with the released package. We can say why rather than treat it as a coincidence: parallel bit-flipping never terminates early on a failure, so every failed trajectory runs the full budget, $\lvert\mathrm{flips}\rvert=T=20$, and both windows equal five iterations on every sample---on this data the two classifiers are the same function. The mode counts therefore no longer carry an exploratory caveat, but they are confirmed \emph{by construction rather than by new evidence}, and we state it that way. The directly measured first-round flip accounting remains separate diagnostic evidence.

\onefig{fig3_mechanism.pdf}
{First-round mechanism. (a) Mean first-round flip accounting per decoding instance (RAND); FRO-2 reduces incorrect flips about $9.5\times$ while reducing correct flips only about $2.7\times$, so the post-round-1 residual is larger. Counts are means over keys and error samples; a single decoding instance still flips an integer number of positions. (b) Failure-mode composition normalized within failures, classified under the preregistered $\lceil T/4\rceil$ window; the 30 mode cells are identical to the archived pre-fix classifier because parallel bit-flipping always exhausts its budget on a failure.}
{fig:mechanism}{0.4em}

\subsection{Transfer and sensitivity}\label{sec:transfer}

Two independent experiments are reported at this same $(r,d,t)=(3001,37,66)$ point, and they must not be conflated. This subsection uses the preregistered transfer run: 30 keys per condition with 3{,}000 errors each. Table~\ref{tab:temporal} instead uses the preregistered expanded confirmation at the same $(r,d)$: 50 regenerated keys with 1{,}000 errors each, which also covers $P_3$. The two therefore differ in keys, error pools and sample size, and their C0 DFRs differ correspondingly ($2.278\times10^{-3}$ against $2.56\times10^{-3}$ in RAND); neither is a subset of the other.

The frozen $+2$ offset transfers to $(3001,37,66)$ without retuning (Fig.~\ref{fig:scope}a). RAND DFR declines from $2.278\times10^{-3}$ to $1.167\times10^{-3}$, TYPE-I from $3.000\times10^{-3}$ to $9.333\times10^{-4}$, and NCW from $1.200\times10^{-3}$ to $4.778\times10^{-4}$; the risk ratios are 1.95 [1.58, 2.42], 3.21 [2.57, 4.14], and 2.51 [1.85, 3.46]. The worst-group log-upper-bound improvement is 0.384, smaller than at the primary point but directionally consistent and statistically resolved.

Iteration-budget sensitivity separates the effect from simple budget exhaustion (Fig.~\ref{fig:scope}c): at budgets 7, 10, and 20 the risk-ratio intervals for RAND and TYPE-I all exclude 1, with point estimates of 1.72--3.17. The budget-20 intervals are wide because several bootstrap resamples contain zero FRO-2 failures; this sparsity is reported rather than suppressed. A fixed-uniform-error key-family analysis likewise favours FRO-2 for random, Type-I, and near-codeword-admissible key batches, with risk ratios of 4.17, 4.44, and 4.33. The near-codeword-admissible batch has only 1.04 times the baseline DFR of random keys under uniform errors, so admission alone does not create the shallow-control effect.

Under the simulation-privileged dynamic-distance semantics, FRO-2 reduces RAND DFR from $7.38\times10^{-4}$ to $4.63\times10^{-4}$ and TYPE-I from $8.53\times10^{-4}$ to $4.25\times10^{-4}$. The 1.59- and 2.01-fold effects are smaller than under static semantics but well resolved ($z=7.94$ and $11.98$). NCW moves from $1.07\times10^{-4}$ to $8.8\times10^{-5}$ and is not statistically resolved ($z=1.36$). The main conclusion therefore does not depend on the static carrier being weaker, although effect magnitude does depend on semantics.

\dblfig{fig5_scope_and_boundary.pdf}
{Transfer and boundary. (a) Key-cluster 95\% intervals at the second operating point $P_2$; the offset was not retuned. (b) Capture gradient in the overlap $\kappa$; the vertical dashed line is $\kappa=d=45$. FRO-2 (circles) is beneficial for $\kappa\le d$ and collapses to the no-effect reference once $\kappa>d$, whereas the near-codeword-aware decoder (squares) is exactly inert for $\kappa\le d$ and becomes decisive at $\kappa=46$. (c) Iteration-budget sensitivity: RAND and TYPE-I exclude 1 at every tested budget. No NCW budget-ablation data were generated, so no third curve is shown.}
{fig:scope}

The effect is not tied to a single operating point. Sweeping the target weight at $P_1'$ on the same code, FRO-2 lowers DFR at every point where failures can be resolved, and the benefit shrinks monotonically as the decoder degrades: risk ratios of 7.00 [3.64, 23.5] at $t=92$, 2.82 [2.28, 3.60] at $t=94$, and 1.87 [1.72, 2.04] at $t=96$, while at $t\le90$ the earlier 30-key sweep saw at most 1 failure in 30{,}000 instances and no contrast was resolvable there. The low-failure-rate end is now resolved on the primary pool itself; the measurements discussed here and in Section~\ref{sec:production} are collected in Table~\ref{tab:added}. Extending the code size, at $(r,d,t)=(6007,56,100)$ --- the largest code in this study and at a lower $t/r$ than any calibrated point --- FRO-2 reduces DFR from $8.13\times10^{-3}$ to $3.38\times10^{-3}$, risk ratio 2.41 [2.09, 2.82]. We also report the full sweep rather than a selected part of it: four further $(r,d)$ points across twelve (point, $t$) cells were attempted, six cells fell below the measurable band (zero C0 failures in 30{,}000 draws) and six above it. Only the three cells at 13--54\% C0 failure produced a resolvable contrast, and there the benefit weakened monotonically toward unity as the decoder degraded (risk ratios 1.77, 1.08 and 1.01 at $t=98$, $102$ and $106$ on the same code). These are knife-edge operating points and the measurable window is narrow; the misses are part of the result.

\subsection{A lower failure-rate point on the same pool}\label{sec:lowdfr}

The sweep above leaves the low-failure-rate direction open: the earlier $t$-sweep found no resolvable contrast at $t\le90$. We therefore repeated $t=90$ on the \emph{same} 92-key pool and the same instances used for the headline result, under a preregistered stopping rule (sample until 120 baseline failures or 100{,}000 errors per key, whichever comes first); the rule triggered at 130 baseline failures in 1{,}840{,}000 paired trials, i.e.\ 20{,}000 errors per key. The baseline DFR falls to $7.07\times10^{-5}$---fourteen times below the primary point---and the risk ratio is 3.71 [2.69, 5.57], with 60 keys improving, 6 worsening and 26 tied (sign test $p=2.7\times10^{-12}$). The interval overlaps the primary point's 4.15 [3.29, 5.38] substantially and the point estimates differ by about $11\%$, so within the measured range the effect is stable in the failure rate rather than strengthening as failures become rarer; the intervals are too wide to exclude a moderate change, and we do not claim one. A cross-pool sweep point at $t=92$ (7.00 [3.64, 23.5], a different 30-key pool) has a point estimate outside the primary point's interval even though the two intervals overlap: pool-to-pool variation of the point estimate is larger than the variation induced by the error weight over this range, which is itself consistent with the key-conditional framing of this paper.

\subsection{A preregistered offset-selection rule validated at a held-out point}\label{sec:rule}

The results so far invite a specific objection: that $+2$ is an empirical tune for a calibrated point rather than the output of a procedure. To answer it we fixed a selection rule, froze its full definition in the released protocol file before the run it covers, and applied it once at an operating point $P_4=(3607,39,76)$ that has never been used for any other purpose in this study. The rule runs no decoding experiment at $P_4$; it is computed from first-round UPC statistics only.

Let $\tau_1$ be the CLASSIC first-round threshold and $q_c(\tau)$ the tail mass of first-round UPC values at $\tau$, separately for true-error positions ($c=1$) and non-error positions ($c=0$). Raising the threshold by $\delta$ avoids $(n-t)\,[q_0(\tau_1)-q_0(\tau_1+\delta)]$ incorrect flips per instance at the cost of deferring $t\,[q_1(\tau_1)-q_1(\tau_1+\delta)]$ correct ones, so we score the exchange rate itself,
\begin{equation}
\begin{aligned}
\rho(\delta)&=\frac{(n-t)\,[q_0(\tau_1)-q_0(\tau_1+\delta)]}
                   {t\,[q_1(\tau_1)-q_1(\tau_1+\delta)]},\\
\delta^{*}&=\arg\max_{\delta\in\{1,2,3,4\}}\rho(\delta),
\end{aligned}
\end{equation}
with ties broken to the smaller offset. The criterion contains no free numerical parameter. An earlier formulation maximizing the asymmetry ratio $A_1/A_0$ per unit $\delta$ was discarded before the real run, because $A_0\to0$ makes it diverge and always select the largest offset.

At $P_4$ the rule selected $\delta^{*}=1$ (Fig.~\ref{fig:rule}a). We then validated it on 48 freshly generated keys with 1{,}000 errors each, with C0, the rule-selected FRO-1, and the frozen FRO-2 decoding identical $(H,e,s)$ instances (Fig.~\ref{fig:rule}b). C0 fails 40/48{,}000 (DFR $8.33\times10^{-4}$). FRO-1 and FRO-2 each fail 8/48{,}000 (DFR $1.67\times10^{-4}$), giving risk ratios of 5.00 [2.64, 16.67] and 5.00 [2.91, 12.00], both excluding one, with paired sign tests at $p=5.6\times10^{-6}$ and $3.0\times10^{-6}$.

Two readings of this result must be separated. What it \emph{supports} is that the first-round effect replicates at a work point that played no part in the study's development, and that a procedure computed from first-round UPC statistics alone independently arrives at a working offset. What it does \emph{not} support is that the rule discriminates well among offsets: at $P_4$, $\rho$ decreases monotonically over $\{1,2,3,4\}$ (0.628, 0.486, 0.409, 0.373), so the criterion selects the smallest offset in its range, and the two arms are indistinguishable at this sample size. The profile shows that the \emph{cheapest} offset already captures the best available exchange rate---not that $+1$ is uniquely optimal. This is consistent with the position taken throughout: FRO-2 is a validated representative intervention, and the best magnitude depends on the operating point.

\onefig{fig6_offset_rule.pdf}
{Prespecified offset rule and its held-out validation. (a) The exchange rate $\rho$ at $P_4$, computed from first-round UPC statistics on 12 development keys without running any decoder; the rule selects $\delta^{*}=1$. The profile decreases monotonically, so the criterion identifies the cheapest sufficient offset rather than discriminating among large ones. (b) Validation on 48 fresh keys $\times$ 1{,}000 errors at $P_4$, with all three arms decoding identical instances. Bars are observed DFR and vertical strokes are one-sided 95\% Wilson caps; labels give failures/trials and key-cluster bootstrap risk ratios against C0.}
{fig:rule}{0.4em}

\subsection{Negative result: round-one statistics do not predict other offsets}\label{sec:counterfactual}

Because the round-one flip counts and the post-round-one syndrome weight can be computed from the initial state and a candidate offset without running that offset's decoder, they invite a counterfactual use: fit $P(\text{fail}\mid\text{round-one statistics})$ on the baseline arm and predict the other offsets from it. On the primary pool such a model classifies baseline failures almost perfectly within its own arm (held-out-key AUC 0.995), yet it predicts the wrong sign for every beneficial offset: predicted risk ratios of 0.79, 0.16 and 0.04 against measured 3.46, 3.86 and 3.97 for $+1$, $+2$ and $+3$ on the same held-out keys. Conditioning on the offset as well, with keys and offsets both held out, still fails (2.18, 1.86 and 0.83 against the same measured values, with the largest offset predicted in the wrong direction), and applying that model to three further configurations --- $t=94$ and $t=96$ on the same code and the larger $(6007,56,100)$ code --- predicts 0.21, 0.27 and 0.33 against measured 2.82, 1.87 and 2.41, outside the measured interval in every case. Section~\ref{sec:discussion} diagnoses the failure and bounds what it implies.

\subsection{Boundary: capture-prone high-overlap stress}\label{sec:boundary}

The shallow NCW condition should not be conflated with this stress test. We fixed one near-codeword-admissible key batch and varied the error construction over uniform errors and $\kappa\in\{12,46,60\}$ (Fig.~\ref{fig:scope}b). FRO-2 improves the uniform and $\kappa=12$ conditions by 3.54 [2.41, 6.00] and 4.75 [2.00, 19.0]. At $\kappa=46>d$, C0 fails on 85.31\% of samples and FRO-2 on 86.50\%; the risk ratio is 0.986 [0.981, 0.991], favouring C0 slightly. At $\kappa=60$ both schedules fail on all 40{,}000 samples.

This is a sharp empirical scope boundary. FRO-2 stabilizes early threshold decisions in the measured recoverable conditions, whereas the high-overlap stress conditions remain highly failure-prone. Methods designed for error-floor or capture recovery address a different failure mechanism \cite{baldelli2026ncw,santini2026flip}. We repeat that these observations show high failure rates under the tested constructions; they do not independently verify a dynamical capture basin.

A structurally different decoder shows that this boundary belongs to the intervention rather than to the experiment. We implemented Algorithm~2 of Baldelli \emph{et al.}~\cite{baldelli2026ncw} --- bit-flipping that recognises when the syndrome has reached a near-codeword and flips that near-codeword's whole support --- under the same CLASSIC threshold, and ran it on identical instances. At $P_1'$, where $\kappa<d$ and capture is impossible by construction, the method is \emph{completely inert}: its near-codeword step never fires across all 834{,}000 instances, its failures reproduce C0 to the last count (282 / 372 / 141), and its mean iterations are identical (4.312, 4.511, 4.206). Under the capture conditions the contrast is sharp (Fig.~\ref{fig:scope}b). At $\kappa=46>d$, C0 fails on 85.31\% of instances while the near-codeword-aware decoder fails on 11.07\%, achieved by a step that fired on 74.2\% of instances. At $\kappa=60$ it also collapses (0.76\% success), so even a purpose-built structural method has a boundary. Across every cell the zero-syndrome success count equals the exact-recovery count, so no miscorrection was observed --- a statement the recovery audit makes checkable rather than assumed.

This dissociation is the clearest support for the framing of this paper. FRO-2 lowers conditional failure where early threshold decisions are recoverable, and is nearly neutral at $\kappa=46$ (risk ratio 0.986); a method built for trapping sets is inert exactly where FRO-2 works and decisive where FRO-2 does not. The two interventions address different failure mechanisms, and neither subsumes the other.

\subsection{Negative control: threshold randomization}\label{sec:negative}

Coin-driven threshold perturbations did not improve conditional risk. On a marginal fixed-threshold carrier --- BF30, i.e. plain parallel bit-flipping with a fixed threshold $b=30$, at the screening operating point $P_1=(4801,45,84)$ --- symmetric jitter $\{-1,0,+1\}$ raised DFR from 4.29\% to 19.67\%, and positive-only jitter $\{0,+1,+2\}$ raised it to 50.5--51.8\% across the three conditions: lower thresholds caused over-flipping and divergence, higher thresholds caused under-flipping and stagnation, and the safe band was narrower than the perturbation support.

On CLASSIC at $P_1'$ the same perturbation pointed the same way but the probe is too small to resolve it: 2/2{,}400 against 7/2{,}400 in RAND, 3/1{,}800 against 7/1{,}800 in TYPE-I, and 6/2{,}400 against 10/2{,}400 in NCW, none significant under Fisher's exact test (Holm-adjusted $p\ge 0.54$). We therefore report the CLASSIC randomization result as \emph{unresolved} rather than as evidence, and rest the negative control on the fixed carrier, where the effect spans an order of magnitude. In formal candidate search, neither coin-selected variant survived. These results reject the randomization hypothesis within the tested design space, but they do \emph{not} prove that every randomized decoder is ineffective---in-place and stateful randomized methods differ substantially from the unstructured threshold mixing tested here \cite{baldi2021inplace}. What follows is that unstructured threshold mixing is not a substitute for a well-positioned deterministic schedule.

\subsection{Positioning against modern decoders}\label{sec:modern}

Table~\ref{tab:modern} partitions the evidence into three layers rather than forming a single ranking: (i) a matched static-semantic comparison of C0, FRO-2, and specification-formula BIKE-Flip on the full formal pool; (ii) a native production reference using official Level-1 constants; and (iii) a semantic reference for Backflip under the C implementation's dynamic-distance rule.

\begin{table*}[t]
\centering
\caption{Modern-decoder comparisons and cross-implementation validation. DFR values are observed; zero-failure cells give rule-of-three upper bounds.}
\label{tab:modern}
\footnotesize
\setlength{\tabcolsep}{4pt}
\resizebox{\textwidth}{!}{%
\begin{tabular}{@{}p{2.3cm}p{2.2cm}p{2.7cm}rp{1.4cm}p{4.4cm}@{}}
\toprule
Layer & Condition\slash parameter & Decoder(s) & Trials & Failures & Reading \\
\midrule
Cross-implementation & $P_1'$, dynamic & Python/C CLASSIC & 20{,}000 / 200{,}000 & 11 / 136 & $z=-0.68$, consistent \\
\midrule
Matched, static & RAND & C0 / FRO-2 / BIKE-Flip / BGF-C3 & 276{,}000 per arm & 282 / 68 / 83 / 287 & FRO-2 lowest; BGF-C3 n.s. vs C0 \\
Matched, static & TYPE-I & C0 / FRO-2 / BIKE-Flip / BGF-C3 & 282{,}000 per arm & 372 / 68 / 94 / 391 & FRO-2 lowest; BGF-C3 n.s. vs C0 \\
Matched, static & NCW & C0 / FRO-2 / BIKE-Flip / BGF-C3 & 276{,}000 per arm & 141 / 31 / 58 / 84 & BGF-C3 beats C0 but trails FRO-2 \\
\midrule
Native Level-1 & all three & spec-exact BIKE-Flip & 8{,}000 per cell & 0 / 0 / 0 & $<3.75\times10^{-4}$; 16 keys/cell \\
Native Level-1 & all three & C0 / FRO-2 & 8{,}000 per cell & 0 / 0 & unresolved vs BIKE-Flip \\
\midrule
Production hold-out & offsets 0, +2, +4 & C0 carrier & $10^6$ per cell & all 0 & bound $<3\times10^{-6}$ \\
Semantic reference & three conditions & Backflip & $10^6$ per cell & 0 & stronger, with more state and cost \\
\bottomrule
\end{tabular}}
\end{table*}

With official Level-1 constants, spec-exact BIKE-Flip also records zero failures in 8{,}000 trials per condition and converges in essentially three iterations. C0 averages 2.05--2.12 iterations and FRO-2 2.40--2.56. BIKE-Flip's stable three-iteration behaviour across conditions illustrates the value of a fully tuned schedule. At the scaled point $P_1'$, where no official constants exist, specification-derived coefficients yield DFRs of $3.007\times10^{-4}$, $3.333\times10^{-4}$, and $2.101\times10^{-4}$: FRO-2 is lower in TYPE-I and NCW and statistically tied in RAND (McNemar $p=0.228$), but this compares a \emph{searched} offset against an \emph{out-of-tuning} derived comparator, and is not evidence of dominance over BIKE-Flip.

At the same instances and the same budget we also ran the GRAY\_BGF structure driven by the parameter-aware CLASSIC threshold (BGF-C3). This is the study's fair structural reference: unlike the specification-exact variant it uses a schedule appropriate to $(r,d)=(4801,45)$. It is statistically indistinguishable from C0 in RAND and TYPE-I (risk ratios 0.983 [0.836, 1.156] and 0.951 [0.837, 1.086]) and improves on C0 in NCW (1.679 [1.362, 2.111]). FRO-2 is lower than BGF-C3 in all three conditions, and BGF-C3 costs about 22\% more iterations (5.28, 5.46, 5.16 against 4.31, 4.51, 4.21). We also record why the specification-exact affine variant is \emph{not} used as a comparator at this point: its published coefficients are BIKE Level-1 values, and the project's own screening record states that no tuned constants exist for $d=45$. Applied here on a reduced sample it fails on 38--60\% of instances, which measures the constant mismatch rather than the algorithm; we report that boundary instead of presenting it as a baseline.

Under reference dynamic semantics, Backflip records zero failures in one million instances per condition, giving upper bounds below $3\times10^{-6}$ and outperforming both C0 and FRO-2 at the scaled point, at the cost of additional state and higher iteration cost. This comparison positions FRO-2 accurately: it is a minimal carrier modification, not the strongest tested decoder.

\subsection{Measurement at the production block parameters}\label{sec:production}

At native BIKE Level-1, the official BIKE-Flip implementation and the C0/FRO-2 carrier arms recorded zero failures in the native comparison (8{,}000 paired instances per condition) and the nine-million-instance hold-out likewise observed zero failures, giving only the per-cell rule-of-three bound below $3\times10^{-6}$; that hold-out was run in the C reference under dynamic-distance semantics, so its bound describes that simulation-privileged carrier rather than the static one, and the convergence figures below carry the same caveat. The measurable production signal is convergence \emph{cost}: two-round convergence falls from 96.0\% at offset $0$ to 68.4\% at $+2$ and 2.2\% at $+4$ for RAND, from 91.7\% to 60.2\% and 1.5\% for TYPE-I, and from 93.6\% to 62.3\% and 2.0\% for NCW (Fig.~\ref{fig:modern}b).

The block parameters themselves can be brought into the measurable regime by raising only the error weight. At $(r,d)=(12323,71)$ we calibrated the baseline failure rate against the error weight on twelve keys with 500 errors each: the failure rate is $1.67\times10^{-4}$ at $t=170$, $2.17\times10^{-3}$ at $t=172$, $1.5\times10^{-2}$ at $t=175$ and $13.9\%$ at $t=180$ (a coarser earlier scan adds 99.9\% at $t=195$). The decodability transition at these block parameters is thus confined to a window of roughly $15\%$ in $t$, and the $10^{-3}$ level sits at $t^{\ast}=172$, about $28\%$ above the production weight of 134. This is a quantitative statement of how far the production operating point sits from the decoding edge, and it explains why no failure is measurable there rather than merely recording that none was seen. At $t^{\ast}$ the first round is two-sided under the production block parameters: on 32 keys with 1{,}500 paired errors each, the $+2$ offset gives a risk ratio of 4.48 [2.97, 7.25] and a single $-1$ gives 0.162 [0.119, 0.209] against baseline 94 failures in 48{,}000. The archived production runs had covered offsets $0$, $+2$ and $+4$ only; the negative direction was measured here for the first time, and it is decisive at the error weight where failures can be resolved.

At the production error weight itself, no failure is measurable for any arm (32{,}000 paired trials each, rule-of-three bound $9.4\times10^{-5}$). What is measurable there is the round-one decision margin: the baseline flips 101.2 positions on average, 87.5 of them true errors and 13.8 false flips; $+2$ reduces false flips to 2.3 (a factor of six) while deferring 25.5 true flips, and $-1$ raises false flips to 31.2. The exchange rate --- avoided false flips per deferred true flip --- is 0.450 at $+2$ against 0.546 at the calibrated point: the same structure, slightly steeper. The measurable production cost is convergence: $+2$ needs 2.446 iterations on average against 2.054 for the baseline, and $-1$ also costs (2.232), so the baseline schedule is the fastest of the three. The specification's own first-round margin behaves the same way under its own decoder: removing it raises the two-round convergence fraction from 0.09\% to 37.35\%, doubling it lowers that fraction further, and applying it in round~1 alone gives 6.14\%. All four ablations record zero failures in 8{,}000 trials each, so at Level-1 sample sizes a measurable convergence cost coexists with a robustness benefit that is not measurable; we report both sides and do not claim the trade is favourable.

\begin{table*}[t]
\centering
\caption{The measurements added for this study: the production block parameters at a raised error weight, the production error weight itself, and the low-failure-rate end on the primary pool. Every arm decodes the same $(H,e,s)$ instances as its baseline; intervals are key-cluster bootstrap 95\% intervals over keys.}
\label{tab:added}
\footnotesize
\begin{tabular}{@{}llrlrl@{}}
\toprule
Setting & Arm & Baseline DFR & Arm failures/trials & RR vs C0 [95\% CI] & Note \\
\midrule
$(4801,45,90)$, 92 keys $\times$ 20{,}000 & $+2$ & $7.07\times10^{-5}$ & 35 / 1{,}840{,}000 & 3.71 [2.69, 5.57] & 60/6/26 keys, $p=2.7\times10^{-12}$ \\
$(12323,71,172)$, 32 keys $\times$ 1{,}500 & $+2$ & $1.96\times10^{-3}$ & 21 / 48{,}000 & 4.48 [2.97, 7.25] & production blocks \\
$(12323,71,172)$, 32 keys $\times$ 1{,}500 & $-1$ & $1.96\times10^{-3}$ & 582 / 48{,}000 & 0.162 [0.119, 0.209] & first negative-direction test \\
$(12323,71,134)$, 32 keys $\times$ 1{,}000 & $+2$ & 0 / 32{,}000 & 0 / 32{,}000 & not measurable & bound $9.4\times10^{-5}$; cost only \\
$(12323,71,134)$, 32 keys $\times$ 1{,}000 & $-1$ & 0 / 32{,}000 & 0 / 32{,}000 & not measurable & bound $9.4\times10^{-5}$; cost only \\
\bottomrule
\end{tabular}
\end{table*}

\onefig{fig7_modern_decoders.pdf}
{Modern-decoder comparison and the production boundary. (a) DFR at the calibrated point $P_1'$ for C0, FRO-2, and specification-formula BIKE-Flip; FRO-2 and BIKE-Flip are unresolved in RAND (McNemar $p=0.23$). (b) Two-round convergence at native BIKE Level-1 as a function of the first-round offset: one million runs per cell with zero failures, so the only measurable signal at that scale is convergence cost rather than DFR.}
{fig:modern}{0.4em}

\subsection{Engineering cost of bounded recovery}\label{sec:engineering}

To separate a minimal schedule intervention from a decoder with additional state, we ran a frozen engineering comparison of C0, FRO-2, STLR, SR-NCW, and BGF-C3 (Table~\ref{tab:eng}) (Fig.~\ref{fig:eng}). STLR (stagnation-triggered local recovery) probes the four largest UPC positions after the first stagnant round; SR-NCW (support-restricted near-codeword recovery) uses the same trigger and budget but restricts candidates to a strictly better-than-canonical near-codeword support in the third condition. BGF-C3 is the GRAY\_BGF structure driven by the parameter-aware CLASSIC threshold (Section~\ref{sec:modern}). The recovery UPC vector is recomputed after the normal round and counted as an additional UPC call. Neither recovery arm receives the true error support.

\begin{table*}[t]
\centering
\caption{Engineering comparison at two calibrated points. Every reported cell has nonzero failures; cost values are Python reference measurements and are not a hardware-independent speed claim.}
\label{tab:eng}
\footnotesize
\begin{tabular}{@{}llrrrrr@{}}
\toprule
Point & Condition & C0 & FRO-2 & STLR & SR-NCW & BGF-C3 \\
\midrule
$P_2$ & RAND        & 35/15{,}000 & 14/15{,}000 & 12/15{,}000 & 12/15{,}000 & 34/15{,}000 \\
$P_2$ & TYPE-I      & 36/15{,}000 & 24/15{,}000 & 23/15{,}000 & 23/15{,}000 & 64/15{,}000 \\
$P_2$ & NCW-shallow & 19/15{,}000 & 9/15{,}000  & 8/15{,}000  & 8/15{,}000  & 13/15{,}000 \\
$P_3$ & RAND        & 19/15{,}000 & 13/15{,}000 & 11/15{,}000 & 11/15{,}000 & 22/15{,}000 \\
$P_3$ & TYPE-I      & 25/15{,}000 & 9/15{,}000  & 9/15{,}000  & 9/15{,}000  & 24/15{,}000 \\
$P_3$ & NCW-shallow & 7/15{,}000  & 1/15{,}000  & 1/15{,}000  & 1/15{,}000  & 4/15{,}000 \\
\bottomrule
\end{tabular}
\end{table*}

STLR and SR-NCW are not merely close. Their failure counts are identical in all six cells (12/12, 23/23, 8/8, 11/11, 9/9, 1/1) and so are their probe counts and trigger rates; the mean iteration count differs in one cell only, by $7\times10^{-4}$ (P2/NCW). The near-codeword support restriction therefore changed no decision on any of the 90{,}000 paired instances, and SR-NCW is effectively a no-op relative to STLR rather than an independent arm; we report it as such. BGF-C3 uses 5.07--5.62 UPC calls per decoding versus 4.1--5.0 for the CLASSIC-based arms, and is not uniformly better than C0. One property of these intervals must be stated. When a bootstrap replicate draws a key set in which the second arm has zero failures, its risk ratio is infinite; the analysis code used for these cells discards such replicates rather than reporting them, and in the $P_3$/NCW cell 36.8\% of replicates are discarded. We therefore re-derived every engineering interval from the archived per-key counts with infinite replicates retained, which makes the affected intervals one-sided. The reading of this paragraph is unchanged --- FRO-2/STLR and FRO-2/SR-NCW still include one in all six cells --- but the finite-only upper endpoints understate the uncertainty in those cells, and the reader should treat them as extending to infinity.

The pairwise intervals are decisive on the question of added value: for FRO-2/STLR and FRO-2/SR-NCW, \emph{none} of the six cell intervals resolves an additional benefit (0/6), so neither stagnation-triggered local recovery nor the near-codeword support restriction yields a resolvable gain under this frozen setup. C0/BGF-C3 resolves in 1/6 cells and shows no consistent DFR advantage. STLR and SR-NCW incur approximately 0.002--0.02 local probes per decoding with trigger rates of 0.05\%--0.51\%; their data-dependent trigger and candidate selection preclude a constant-time claim.

The engineering interpretation is deliberately bounded: first-round conservatism remains the most stable intervention with the smallest structural change; stagnation-triggered recovery changes the DFR/cost trade-off in selected cells; the near-codeword support restriction does not yield a consistent gain. BGF-C3 demonstrates the cost of adding black/gray substeps under the same static threshold but is not an official tuned BGF result.

Finally, we measured the two carriers in the project's specification C reference, compiled with \texttt{zig cc -O2} for \texttt{x86\_64-windows-gnu} with AVX disabled, single-threaded, 500{,}000 decodes per repetition and the minimum wall time over three repetitions on an otherwise idle machine. C0 decodes at 7{,}134 per second and FRO-2 at 6{,}563, a throughput ratio of 1.087; mean iterations are 4.191 against 4.471 (ratio 1.067) and the P99 is six for both. A compiled implementation therefore pays roughly 7--9\% for the intervention, which agrees with the Python-level iteration count and shows that the cost is a genuine convergence effect rather than interpreter overhead. Two caveats bound this measurement. The reference uses dynamic-distance threshold semantics, in which the threshold may consult the current distance to the true error; that is simulation-privileged and is not the static semantics of the main experiment, so only the C0-versus-FRO-2 ratio is interpreted, not the absolute DFR. And the GRAY\_BGF build available in the same binary uses the mis-scaled affine constants discussed above, so its 13.2 mean iterations and 47.9\% budget-exhaustion rate measure that mismatch rather than BGF's cost; we report the number only to make the point explicit.

\onefig{fig8_engineering_tradeoff.pdf}
{Engineering cost--DFR trade-off. (a) Each point is one method in one operating-point and condition cell: abscissa is UPC passes per decode, ordinate is DFR. (b) Mean risk ratio over the six cells for selected method pairs, with 95\% intervals; the right-hand label gives the number of cells whose interval excludes one. For FRO-2 against STLR and SR-NCW, none of the six intervals resolves an additional benefit.}
{fig:eng}{0.4em}

% ===========================================================================
\section{Discussion}\label{sec:discussion}

\subsection{What the evidence supports}

Across the primary formal pool, the pure-fresh subset, a second calibrated point, two decoder semantics, multiple iteration budgets, and three key families, a first-round $+2$ offset generally lowers the conditional failure rate of the tested CLASSIC carrier. Every key-cluster interval in the primary formal experiment and in the prespecified $P_2$ transfer experiment exceeds a C0/FRO-2 risk ratio of 1, and most keys improve individually; the dynamic-semantics NCW condition and several sparse engineering cells remain statistically unresolved.

More important is how that effect decomposes. The uniform offset is not harmful because it touches round~1 --- an arm that offsets only round~2 onwards, leaving round~1 untouched, is equally destructive --- but because it displaces every later round, which prevents the deferred corrections from being recovered. Symmetrically, the first round is a two-sided window: a permissive first decision is as damaging as a conservative one is helpful. The mechanistic proposition is therefore that a \emph{single}, well-placed early decision changes the state entering later iterations, and that the later schedule cannot absorb a sustained displacement.

The mechanism measurements refine that interpretation. FRO-2 is conservative in the strict sense: it makes fewer total first-round flips, far fewer wrong flips, and leaves more real errors behind. A decoder can therefore improve its final success while looking worse on immediate residual weight. This cautions against using early error reduction alone as a schedule objective; an appropriate objective must account for the quality of the residual support and the ability of later rounds to recover.

A counterfactual use of round-one statistics fails, and the reason matters. Because the round-one flip counts and the post-round-one syndrome weight are computable from the initial state and a candidate offset without running that offset's decoder, it is natural to fit $P(\text{fail}\mid\text{round-one statistics})$ on the baseline arm and predict other offsets from it. On the primary pool such a model classifies baseline failures almost perfectly within its own arm (held-out-key AUC 0.995) yet predicts the wrong sign for every beneficial offset --- predicted risk ratios of 0.79, 0.16 and 0.04 against measured 3.46, 3.86 and 3.97 for $+1$, $+2$ and $+3$ on the same held-out keys. Conditioning on the offset as well, with keys and offsets both held out, still fails (2.18, 1.86 and 0.83 against the same measured values, with the largest offset predicted in the wrong direction), and the same model applied to three further archived configurations --- $t=94$ and $t=96$ on the same code and the larger $(6007,56,100)$ code --- predicts 0.21, 0.27 and 0.33 against measured 2.82, 1.87 and 2.41, outside the measured interval in every case. The diagnosis is a post-treatment conditioning problem: a low count of flipped true errors means ``hard instance'' under the baseline, but ``deferred by the intervention'' under a raised threshold, so the conditional failure probability does not transport across offsets. This is an empirical instance of the hazard \cite{annechini2025closedform} warns about, and it bounds what round-one statistics can support: the aggregate exchange rate explains why a conservative first round can help and tracks the sign across every measured point, but no validated quantitative predictor was obtained and none is claimed.

One objection deserves an explicit answer. The production decoder sustains its margin $\delta$ across all iterations, whereas we report that sustaining an offset destroys the decoder. Both statements concern perturbations of \emph{particular} schedules. Our sustained-offset arm displaces the later rounds of a schedule whose thresholds are derived from the current syndrome weight through the binomial model: every later round of that carrier is individually calibrated, and shifting all of them breaks the calibration. The specification's schedule has a different shape --- its later iterations are floored at $(v{+}1)/2$, its margin is present from the first iteration by construction, and its calibration target is the whole schedule rather than each round separately. The correct reading is therefore that a \emph{mid-flight} displacement of a calibrated per-round schedule is destructive, not that any schedule containing a constant margin must fail.

The negative control is equally informative. Random threshold choice changes the conditional mixture of trajectories but does not automatically reduce key dependence or worst-group risk; when the viable threshold band is narrow, mixing schedules merely samples failure cliffs. Deterministic timing instead exploits a repeatable asymmetry between the initial and later syndrome states.

\subsection{Scope and security relevance}

The experiments measure decoder-layer conditional risk. They do not demonstrate a practical key-recovery attack or an IND-CCA security proof, and they do not negate BIKE's implicit-rejection protections. A DFR change can translate into a system-level reliability or security consequence only when the deployed parameters, rejection semantics, constant-time constraints, and error-feedback interface match the experiment. FRO-2 changes only the first-round threshold calculation and requires no persistent state, making it a minimal algorithmic modification to CLASSIC-like decoders; it nevertheless incurs measurable iteration and early-convergence costs. Where tuned BIKE-Flip or Backflip is available, these experiments provide no evidence that FRO-2 should replace it.

The extrapolation boundary is equally important. The main DFR lies at $10^{-4}$--$10^{-3}$ calibrated pressure points, far above BIKE's cryptographic target, and no production-tail extrapolation is attempted. The tested high-overlap structured-error conditions have very high failure rates for both C0 and FRO-2; these results do not establish a general capture basin, and the intervention is not shown to be an error-floor defence. Production-scale zero-failure runs exclude only degradation above the achieved resolution and cannot demonstrate a benefit. RAND primarily supports the sampled random-key generator; TYPE-I is an adversarial structural probe; NCW-shallow is a shallow-intersection control that changes both key admission and error construction. The fixed-uniform-error near-codeword family supplies only gathering-adjacent structural evidence, not an exact gathering-key result.

\subsection{Limitations}

The primary conclusions rely on two operating points chosen for measurable failures; a third calibrated point is used only for supplemental sensitivity and expanded confirmation. Other limitations include a researcher-defined candidate search space; supplemental cells with only 15--30 keys; a cross-implementation check performed within the same research project rather than by an independent laboratory; and first-round labels that use the true error support and are therefore diagnostic rather than deployable.

Three boundary statements are specific to the added results. First, the production-parameter point is the production \emph{code} at a raised error weight, not the production operating point, and the operating point was selected by a bracketed bisection whose samples are a subset of the main run's instances; the effect measured there is not evidence about the designed-weight failure rate. Second, the production-error-weight statistics are mechanism measurements computed without decoding for the counterfactual offsets, and the convergence-fraction differences at that weight are the only statistically resolved production signals we have. Third, on the primary pool the per-key failure counts are consistent with homogeneous Poisson variation (variance-to-mean 0.98; the 95\% upper bound on the across-key coefficient of variation of the failure rate is about 30\%), and the per-key failure sets of the two arms barely correlate, so ``key-conditional'' here describes the inference, not a demonstrated spread of key-level failure rates.

Future work should add an independently calibrated point or a semi-analytic tail bound; implement a precisely defined gathering-key construction; obtain an external blind replication; report runtime, UPC counts, timing distributions, constant-time constraints, and memory-access behaviour under matched semantics and budgets; and give the local-recovery prototype an explicit structural score that can be shown to track near-codeword structure.

% ===========================================================================
\section{Conclusion}\label{sec:conclusion}

A two-count increase in the first-round threshold, followed by an otherwise unchanged CLASSIC schedule, reduces the measured conditional DFR by approximately 4.1--5.5-fold at the primary calibrated point. The direction persists under key-cluster inference, a pure-fresh analysis, a second calibrated point, and a reference semantics. Direct measurement explains the trade-off: FRO-2 suppresses incorrect first-round flips by about $9.5\times$ while giving up correct flips by only about $2.7\times$, accepts a larger immediate residual, and relies on later rounds for recovery.

The same experiments delimit what the intervention does not do. Threshold randomization is ineffective or harmful in the tested regimes. The high-overlap structured-error stress conditions remain highly failure-prone, and a near-codeword-aware decoder is the appropriate tool there. At the production error weight the available sample size detects no failures for any tested offset, so no benefit is established; what is established there is the round-one decision margin and a convergence cost, and a counterfactual model of the offset effect built from round-one statistics fails out of sample. Tuned BIKE-Flip and stateful Backflip remain the appropriate performance references.

The supported conclusion is not a universal production advantage but a local principle that replicated across every operating point tested here. It has four parts: at measurable CLASSIC operating points, a one-shot conservative offset on the first round can stabilize part of the decoding trajectory; the first-round window is two-sided, so a permissive first decision is consequential in the opposite direction; sustaining an offset across later rounds destroys the decoder whether or not round 1 is touched; and the principle survives at the production block parameters, where the window is two-sided once the error weight is raised into the measurable band, while the failure rate at the production weight itself remains out of reach. The tested high-overlap structured-error conditions remain failure-prone. Key-dependent decoder studies should pair outcomes on identical instances, split at the key level, treat keys as clusters, report zero-failure upper bounds, and test explicit boundaries.

% ===========================================================================
\section*{Data and Code Availability}
The implementation, frozen configurations, seed derivation rules, raw per-key counts, and analysis code are archived at \url{https://doi.org/10.5281/zenodo.22296719}, with a development mirror at \url{https://github.com/Leekj12410/qc-mdpc-experiment}. The archive covers the formal primary and round-control runs, the transfer, expanded-confirmation, mechanism, capture and iteration-ablation experiments, the exact-recovery audit, and the BIKE-Flip, BGF and near-codeword-aware comparisons. This manuscript's own artefacts are released with it: the figure source tables, the script that regenerates them from the archived counts, and the script that checks every quoted value against those tables. Each number in this paper is traceable through that chain, and the checking script is included so the mapping can be re-verified rather than taken on trust.

\section*{Author Contributions}
[Pending before submission: add author names, affiliations, and each author's CRediT roles.]

\section*{Competing Interests}
[Pending before submission: confirm and state the authors' competing-interest declaration.]

\bibliographystyle{IEEEtran}
\bibliography{refs}

\end{document}
